% !TeX root = ../main.tex
\chapter{So sánh và thảo luận kết quả}
\label{chap:comparison}

\section{Bảng so sánh}

Tất cả số liệu trong Bảng~\ref{tab:model-comparison} được sinh tự động từ các
model output và lưu tại \path{results/model_comparison.csv}.

\begin{table}[H]
\centering
\caption{So sánh baseline và ba mô hình cải thiện}
\label{tab:model-comparison}
\resizebox{\textwidth}{!}{%
\begin{tabular}{lcccccccc}
\toprule
Model & Train Acc. & Test Acc. & Error & Weighted F1 & Mal. Prec. & Mal. Recall & Depth & Leaves \\
\midrule
Baseline & 1,0000 & 0,9123 & 0,0877 & 0,9130 & 0,8478 & 0,9286 & 7 & 19 \\
Max Depth & 0,9868 & \textbf{0,9386} & \textbf{0,0614} & \textbf{0,9387} & 0,9070 & \textbf{0,9286} & 4 & 11 \\
Min Samples Leaf & 0,9802 & 0,9123 & 0,0877 & 0,9127 & 0,8636 & 0,9048 & 6 & 13 \\
Pruning & 0,9736 & \textbf{0,9386} & \textbf{0,0614} & 0,9384 & \textbf{0,9268} & 0,9048 & \textbf{3} & \textbf{6} \\
\bottomrule
\end{tabular}}
\end{table}

\begin{figure}[H]
\centering
\includegraphics[width=\textwidth]{model_comparison.png}
\caption{Trade-off giữa accuracy và số leaf của các mô hình}
\label{fig:model-comparison}
\end{figure}

\section{Mô hình tốt nhất}

Max Depth và Pruning cùng đạt test accuracy 93,86\% và error rate 6,14\%.
Notebook không dùng thứ tự dòng để xử lý tie mà áp dụng quy tắc minh bạch: ưu
tiên malignant recall, sau đó weighted F1 và độ đơn giản. Với bài toán chẩn
đoán, bỏ sót malignant là loại lỗi đáng chú ý hơn.

Max Depth được chọn làm mô hình chính vì malignant recall 92,86\%, cao hơn
90,48\% của Pruning, đồng thời weighted F1 nhỉnh hơn (93,87\% so với 93,84\%).
Hai mô hình đều dự đoán sai 7 mẫu, nhưng Max Depth bỏ sót 3 ca malignant và
cảnh báo nhầm 4 ca, trong khi Pruning bỏ sót 4 ca malignant và cảnh báo nhầm
3 ca. Vì vậy, Max Depth phát hiện thêm một ca malignant mà không làm tăng tổng
số lỗi.

\section{Trade-off và ý nghĩa}

Pruning là phương án cạnh tranh mạnh khi khả năng diễn giải hoặc kích thước mô
hình được ưu tiên: chỉ 6 leaf và depth 3, so với 11 leaf và depth 4 của Max
Depth. Pruning cũng có malignant precision cao hơn, nghĩa là tỷ lệ dự đoán
malignant đúng cao hơn trong số các ca bị gắn nhãn malignant. Vì vậy, kết luận
không phải một phương pháp luôn tốt hơn; lựa chọn phụ thuộc chi phí của từng
loại lỗi và yêu cầu về độ đơn giản. Đây cũng là biểu hiện của bias--variance
trade-off: chấp nhận tăng nhẹ bias trên training set để giảm variance và cải
thiện khả năng tổng quát hóa trên dữ liệu chưa quan sát.

Min Samples Leaf cho thấy một ``improvement method'' không nhất thiết phải tăng
accuracy. Giá trị học thuật nằm ở quy trình thử nghiệm, đo lường và giải thích:
phương pháp này giảm cây nhưng không giảm error trên test split hiện tại.
